%!TEX TS-program = PGF
%!TEX encoding = UTF-8
\documentclass{standalone}
\input{tikzstandalone}
\begin{document}

\def\pKeval{14}% produit ionique de l'eau
\def\pKaa{8.4}% Ka1

\def\courbepHCa{0.15}% becher
\def\courbepHVa{10.0} % volume d'acide (mL)
\def\courbepHCb{0.10} % burette
\pgfmathsetmacro{\Veq}{\courbepHCb*\courbepHVa/\courbepHCb}
	
\begin{tikzpicture}[
	declare function={%déclaration de fonctions :
	h(\pH)=10^(-\pH);% h concentration des ions H3O+
	x0(\pH)=1/(1+10^(\pH-\pKaa)); % fraction acide
	x1(\pH)=1/(1+10^(\pKaa-\pH));% fraction base
	w(\pH)=10^(-\pKeval+\pH)-h(\pH); % w-h
	frac_base(\pH)=x1(\pH) + w(\pH)/\courbepHCa;
	V_verse(\pH)=frac_base(\pH)*\Veq;% volume de base versée
}]
	
	\begin{axis}[% mise en forme du graphique :
		xmin=0,xmax=25,
		ymin=0,ymax=14,
		xlabel={volume $V_b$ versé (\si{mL})},
		ylabel={\pH},
		ytick distance=2,
		minor y tick num=1,% nombre de graduations secondaires intermédiaires
		minor x tick num=4,% nombre de graduations secondaires intermédiaires
		grid,
		grid=both,% grilles primaire et secondaire
		title={titrage pH-métrique d'un acide faible},
		axis line shift=3pt,
		xtick align=outside,
		ytick align=outside,
		tickpos=left,
	]
		\addplot[domain=0:14, thick, samples=100] ({V_verse(x)},{x});% volume versé en x, et pH en y
		% équivalence :
		\coordinate (equiv) at (\Veq, 7);
		\draw[violet] (\Veq, 0) -- (equiv);
		\filldraw [violet](equiv) circle (1pt) node[above right, font=\scriptsize \sffamily]{équivalence};
		\draw[red] (0,14) node[anchor=north west, text width=2.7cm, align=left, fill=white, draw, inner sep=1mm] {$V_a=\SI{\courbepHVa}{mL}$\\$C_b=\SI{\courbepHCb}{\M}$};
	\end{axis}
	
\end{tikzpicture}

\end{document}
